\section{Conclusions}
\label{sec:conclusions}

\subsection{Summary of Findings}

This paper demonstrates that machine learning models can effectively predict
Formula 1 race outcomes, substantially outperforming a grid-position baseline.
Our XGBoost model achieves 92\% top-3 accuracy on the 2025 test season,
confirming that structured race dynamics are learnable from historical telemetry
and performance data.

Key findings include:

\begin{itemize}
    \item \textbf{XGBoost outperforms neural networks} on this tabular dataset,
          consistent with the broader literature on gradient boosting for structured
          prediction tasks.
    \item \textbf{Weather features are disproportionately important} in wet races,
          where models with weather features outperform the baseline by 34 percentage
          points in top-3 accuracy.
    \item \textbf{Recent form matters more than career history}: rolling 5-race
          averages rank higher in feature importance than all-time averages.
    \item \textbf{Temporal validation is essential}: naive random splitting
          inflates reported accuracy by 8--12\% by leaking future season patterns
          into training.
\end{itemize}

\subsection{Limitations}

Several limitations constrain the current approach:

\begin{enumerate}
    \item \textit{Strategy decisions}: pit stop timing, undercut/overcut attempts,
          and safety car responses are not modelled, yet critically affect race outcomes.
    \item \textit{Real-time data}: predictions are made pre-race; in-race adaptation
          based on live telemetry is not implemented.
    \item \textit{Data availability}: FastF1 telemetry quality varies by year and
          session; some early-season races in 2020 have incomplete telemetry.
    \item \textit{Rare events}: crashes, mechanical failures, and red flags are
          low-probability events that remain difficult to predict.
\end{enumerate}

\subsection{Future Work}

Promising directions for extending this research include:

\begin{itemize}
    \item \textbf{Real-time prediction}: updating predictions lap-by-lap using
          live telemetry feeds.
    \item \textbf{Strategy simulation}: modelling optimal pit stop windows as a
          separate optimization layer on top of pace predictions.
    \item \textbf{Driver DNA modelling}: characterizing individual driver
          wet-weather preference, overtaking aggression, and tire management
          style as latent embeddings.
    \item \textbf{Ensemble methods}: combining XGBoost and neural network
          predictions via stacking for marginal accuracy gains.
    \item \textbf{Uncertainty quantification}: providing calibrated confidence
          intervals rather than point predictions.
\end{itemize}

The code, data pipeline, and trained models are available in the accompanying
repository, enabling reproducibility and extension by the research community.
